

The widespread adoption of AI has been largely driven by deep learning methods \cite{Lecun2015436}, but deep neural networks (DNNs) have high computational complexity and energy requirements making them unsuitable for energy-constrained systems such as embedded and edge devices. As an alternative, the Tsetlin Machine (TM) offers low computational complexity and high explainability, while keeping competitive performance compared to other learning algorithms \cite{GranmoOle-Christoffer2018TTM-}.

Introduced in 2018 by Ole-Christoffer Granmo, the TM is built on propositional logic and finite-state automata \cite{GranmoOle-Christoffer2018TTM-}. Its operations rely on bit manipulation, making it well-suited for hardware (HW) implementations. Various accelerators have been implemented on FPGA and ASIC platforms \cite{AndersTunheimSvein2025AA86,TunheimSveinAnders2023CTMT,TunheimSveinAnders2025TMIC}, demonstrating the potential for high-speed computations and low power consumption. However, these fixed-function implementations lack flexibility, requiring significant architectural redesign for each new TM variant.

Processor-oriented acceleration through custom instruction extensions offers a flexible alternative that preserves programmability. The RISC-V instruction set architecture (ISA) enables the design of customizable processors that can be extended with domain-specific instructions, and research on RISC-V-based ML accelerators has demonstrated significant performance and energy-efficiency gains \cite{NunesWillianAnaldo2025AMLw, WangShihang2024OCCU,GarofaloAngelo2021XEEE}. 

This project implements a processor-level accelerator for TM inference by designing a custom RISC-V instruction that targets clause evaluation, the most computationally intensive operation in the inference pipeline. The custom instruction is implemented using the NEORV32 processor's Custom Functions Unit (CFU) and evaluated on a Xilinx Artix-7 100T FPGA. By accelerating clause evaluation while maintaining the processor's ability to execute arbitrary code, this approach seeks to balance performance improvements with the flexibility needed to support different TM architectures alongside other computational tasks on edge devices.